% BANKS-SOURCE: pdf=126 print=116 kind=section id=8.7
\subsection{Symmetries and symmetry breaking in the strong interactions}

% BANKS-SOURCE: pdf=126 print=116 kind=prose
The chiral and anomaly formulas in this section are the $D=4$ specialization;
the functional-integral expressions are written in Euclidean signature.

The modern theory of the strong interactions is called quantum chromodynamics
(QCD). It is an SU(3) gauge theory with $N_F$ Dirac quark fields transforming
in the triplet representation of SU(3). The quark mass matrix has the form
% BANKS-SOURCE: pdf=126 print=116 kind=equation id=8.5
\begin{equation}
  \bar q_R^{\,j}M_{ij}q_L^{\,i}+\text{h.c.}
  \label{eq:8.5}
\end{equation}
Here $q_{L,R}^{\,i}=(1\pm\gamma_5)q^i/2$. The mass matrix could be a general
complex matrix if we do not insist on other symmetries. The kinetic terms are
classically invariant under U($N_f$) $\times$ U($N_f$) transformations on the
quarks. Using these we can make the mass matrix diagonal, with positive
entries. However, as we will discuss, the $U_A(1)$ transformation
\[
  q^i\longrightarrow \mathrm{e}^{\mathrm{i}\alpha\gamma_5}q^i
\]
does not leave the functional measure of the quarks invariant. It changes the
Euclidean action by (all traces in this section are taken in the fundamental
representation of SU(3) unless otherwise noted)
\[
  \delta S=(2N_f)\,\mathrm{i}\,\frac{\alpha}{32\pi^2}
  \int\!\operatorname{tr}(F_{\mu\nu}\widetilde F^{\mu\nu}) .
\]

% BANKS-SOURCE: pdf=127 print=117 kind=prose
where $\widetilde F^{\mu\nu}=\tfrac12\epsilon^{\mu\nu\lambda\kappa}
F_{\lambda\kappa}$. We will see in the final chapter of this book that finite
action configurations of the gauge field have
$[1/(32\pi^2)]\int\operatorname{tr}(F_{\mu\nu}\widetilde F^{\mu\nu})$, which is
quantized in integer units, so unless $2N_f\alpha=2\pi n$ for the
transformation which makes the determinant of the quark mass matrix real, we
should expect a term
\[
  \frac{\theta_{\rm QCD}}{32\pi^2}
  \int\!\operatorname{tr}(F_{\mu\nu}\widetilde F^{\mu\nu})
\]
in the QCD action. $\theta$ could be any number between 0 and $2\pi$, but a
dimensional-analysis estimate of the neutron electric dipole moment suggests
that $\theta<10^{-9}$. There have been several proposed explanations of why
$\theta_{\rm QCD}$ should be so small, none of which is completely
satisfactory. This is called the \emph{strong CP problem}.

Apart from the mysterious parameter $\theta_{\rm QCD}$, QCD appears to have a
dimensionless coupling, which appears in
\[
  \frac{1}{4g^2}\operatorname{tr}(F_{\mu\nu}F^{\mu\nu}) .
\]
However, in the course of renormalization of the theory, this will be replaced
by a scale $\Lambda_{\rm QCD}$.\footnote[13]{This mysterious-sounding process
is called \emph{dimensional transmutation}. It is a consequence of the fact
that the dimensions of operators in quantum field theory receive quantum
corrections, as we will see in Chapter 9.} Experiments on high-energy
scattering suggest that this scale is of order 150 MeV. We have already
alluded to the fact that the properties of pions, the lightest hadron, can be
best understood if we assume that two of the quark masses, $m_u$ and $m_d$,
are much smaller than $\Lambda_{\rm QCD}$. Data on K and $\eta$ mesons are
also compatible with the idea that the strange-quark mass can be treated as
small. In the limit in which we set all of these masses to zero, QCD has an
SU(3) $\times$ SU(3) chiral symmetry. Much of low-energy hadron physics can
be understood if we assume that this symmetry is broken spontaneously to the
diagonal vector subgroup SU(3), with the explicit breaking of the latter
symmetry by quark masses treated in first-order perturbation theory.\footnote[14]{One
also has to take into account breaking of the isospin subgroup, SU(2), by
electromagnetism.} The light pseudo-scalar mesons are interpreted as the
pseudo-Nambu--Goldstone bosons of this broken symmetry.

Our purpose in this section is to outline some of the theoretical arguments
which support this conclusion. We will assume that quark confinement occurs in
QCD. The arguments of the last section gave us a plausible picture of why this
is so. More compelling is the experimental fact that, although we can see
evidence for the quark structure of hadrons in experiments at very high energy
and momentum transfer, and in the structure of the hadron spectrum itself, no
experiment has ever produced an isolated quark.

The first part of our argument is that the vector SU(3) currents do not suffer
spontaneous breakdown. To see this, note that if we gave the up, down, and
strange quarks equal mass, $m$, SU(3) would still be a good symmetry. We can
then ask whether this

% BANKS-SOURCE: pdf=128 print=118 kind=prose
symmetry is spontaneously broken, that is, whether there could be a NG boson
pole in the two-point function
\[
  \int\!\dd^4x\,\mathrm{e}^{\mathrm{i}px}
  \langle J_A^\mu(x)J_B^\nu(0)\rangle .
\]
In the functional integral formalism this is calculated as the average of
\[
  \operatorname{Tr}(\gamma^\mu T_A S(x,0)\gamma^\nu T_B S(0,x)),
\]
over gauge field configurations. $S(x,0)$ is the quark propagator in an
external gluon field. The measure of integration over gluons is
\[
  \mathrm{e}^{-\frac{1}{4g^2}\int\operatorname{tr}F^2}
  \det(\mathrm{i}\gamma^\mu D_\mu+\mathrm{i}m),
\]
which is positive (Problem 8.12). Note that we have assumed
$\theta_{\rm QCD}=0$. The result we will obtain is true for at least a small
range of $\theta_{\rm QCD}$ around 0.

The heart of the argument, which is due to Vafa and Witten \cite{banks-ref-99}, is that
$S(x,0)$ falls off exponentially at large $x$, with a bound on the exponent
that is independent of the gauge configuration. This is obvious in
perturbation theory around vanishing gauge field and for gauge fields that fall
off rapidly at infinity. Vafa and Witten proved that it is true in general. It
follows from general properties of probability measures that the averaged
two-point function also falls exponentially, so the two-point function cannot
have a pole at zero momentum. Thus, the vector SU(3) is not spontaneously
broken for any finite $m$. But this means that it cannot be broken for $m=0$,
because turning on $m$ does not break the symmetry, so NG bosons present at
$m=0$ would persist for finite $m$. To prove that the axial symmetries are
spontaneously broken for $m=0$, we will have to make a short detour through
the subject of \emph{anomalies}.

\BanksImplicitHook{I-CH08-010}
